\section{The Bakeoff: Eight Combinations Across Four States}
\label{sec:bakeoff}

\paragraph{Evidence status and relationship to the governed national bakeoff.}
This May 2026 research table is not the later NRS v0.3 national bakeoff. It
mixes confirmed single-profile results with explicitly estimated and pending
cells, uses tract-level political and demographic inputs, and does not provide
a common-block national comparison. Accordingly, cross-configuration patterns
below are exploratory hypotheses wherever they depend on a marked cell. The
governed NRS v0.3 package compares only NRS with official CD118 assignments on
a frozen 2020 block universe and does not validate this eight-configuration
matrix.

\subsection{Preregistered Wisconsin Structure Proof Slice}

On 2026-08-09, before inspecting comparative results, we froze a separate
common-input proof slice on Wisconsin's 2020 congressional tract graph. The
slice holds 1,542 tracts, eight districts, total population, geographic edge
weights, single outer search, and requested seed 0 fixed while varying four
implemented structure families. Its package regenerated exactly after
canonical district relabelling. Table~\ref{tab:wi-structure-proof} reports the
validity gate, not a ranking.

\begin{table}[h]
\centering
\caption{Post-remediation Wisconsin structure validation. The archived initial
  run retained two audit failures; all four repaired executions pass.}
\label{tab:wi-structure-proof}
\small
\begin{tabular}{llll}
\toprule
\textbf{Structure} & \textbf{Final seed} & \textbf{Audit} & \textbf{Failure witness} \\
\midrule
Standard bisection & 2 & pass & none \\
Ratio-optimal & 0 & pass & none \\
Ratio-optimal-area & 0 & pass & none \\
Prime-factor & 0 & pass & none \\
\bottomrule
\end{tabular}
\end{table}

The standard run documents an automatic balance retry from requested seed 0
to accepted seed 2. The standard and ratio-optimal assignments agree on 837 of
1,542 tracts (54.2802\,\%) after maximum-overlap district-label matching.
Across all six post-remediation pairs, agreement ranges from 54.0208\,\% to
66.2776\,\%. This establishes reproducible assignment difference on this
input, not a causal effect or a preferred plan.

The first governed run is preserved separately. In that run,
ratio-optimal-area disconnected a district, while prime-factor disconnected a
district and produced 0.704852\,\% population deviation against the advertised
0.5\,\% tolerance. The prime-factor manifest nevertheless marked population
balance valid. Inspection found that ApportionRegions used a hidden 3\,\%
final validation tolerance, the population flag was hard-coded, and the two
affected METIS paths did not consistently use the contiguity-capable k-way
routine. Final validation now uses the configured tolerance, the flag is
derived from the audit certificate, and k-way contiguity options are enabled.
The repaired four-pass package and its exact regeneration close the one-State
implementation gate; they do not authorize national expansion without a newly
frozen schedule.

The governed evidence is in
\nolinkurl{docs/experiments/neutral-algorithm-family-bakeoff-wi-2020/} under
protocol \texttt{neutral-algorithm-family-bakeoff-v1}. Because this slice uses
a different metric and input contract, it does not replace any estimated or
pending cell in the older tables below.

\subsection{Frozen National Structure Pilot}

A second protocol, \texttt{neutral-algorithm-family-national-bakeoff-v1},
freezes the national multi-district cohort, a staged schedule, failure policy,
and claim boundary before expansion. Its pilot orders RI, NE, CT, KY, SC, WI,
AZ, and CA, spanning two through 52 congressional districts. Each State holds
its 2020 tract graph, seat count, total population, geographic edge weights,
single outer search, and requested seed 0 fixed across the same four structure
families. The complete target is the 44 States with at least two House seats;
the six at-large States are design exclusions rather than passes.

The pilot is an engineering gate selected to exercise increasing scale, prime
and composite seat counts, the Wisconsin anchor, and the maximum-size State.
It is not a probability sample. A passing pilot can establish bounded
multi-State implementation validity and authorize the already frozen full
phase; it cannot estimate national prevalence, rank the structures, or support
geometric, electoral, demographic, VRA, causal, or legal conclusions. Only the
complete 44-State phase could establish coverage of that frozen cohort, and it
still could not identify a generally best algorithm.

The retained initial pilot attempt passed 27 of 32 cells and six of eight
States. It exposed a recursive contiguity defect in equal non-NRS splits: an
early recursive-METIS split could fragment a child region before a later
contiguity-constrained call. After routing all non-NRS splits through k-way
METIS with contiguity and minimum-connectivity options, focused Arizona and
California probes passed and the clean governed matrix passed all eight States
and all 32 cells. Across its 48 within-State structure pairs, maximum-overlap
assignment agreement ranges from 49.0634\,\% to 100\,\%; alternative edge-cut
ratios to standard-bisect range from 0.856479 to 5.441826. These are descriptive
mechanics, not pooled estimates. An independent verifier reran all 32 cells and
passed byte-identical normalized comparison.

The authorized full phase subsequently executed all 44 multi-district States
and passed all 176 cells without a seed retry. Across 264 within-State pairs,
maximum-overlap assignment agreement ranges from 47.6257\,\% to 100\,\%; the
132 alternative edge-cut ratios to standard-bisect range from 0.623301 to
5.441826. An independent execution reran all 176 native cells and regenerated
the normalized evidence exactly. Edge cuts are reported to five decimal places,
while canonical assignments are byte-compared. Thus the package establishes
complete coverage of the frozen cohort. These
descriptive mechanics cannot support a generally best algorithm or political,
geometric, demographic, VRA, causal, or legal conclusions.

\subsection{Experimental Design}

We evaluate eight algorithm combinations on four states that concentrate the
full range of redistricting legal challenges.
Wisconsin (8D, 50.3\,\% D vote) is the canonical competitive state with a
compact Milwaukee metro core.
North Carolina (14D, 49.3\,\% D vote) has the Piedmont corridor geographic
structure, a history of VRA litigation (\emph{Thornburg v.\ Gingles}, 1986),
and an active state-court partisan neutrality standard.
Georgia (14D, 50.1\,\% D vote) has a substantial Black minority (33\,\% of
VAP), significant VRA exposure after the Stacey Abrams-era map litigation,
and rapid Sun Belt population growth.
Alabama (7D, 37.1\,\% D vote) is the \emph{Allen v.\ Milligan} (2023) state:
the Supreme Court required two Black-opportunity districts from a state with
27.5\,\% Black VAP.

\paragraph{Eight combinations.}
\begin{enumerate}
  \item \textbf{Baseline}: Standard balanced bisection, unweighted edges
    (B.1 standard).
  \item \textbf{Geo+GeoSection}: Geographic edge weights + GeoSection ratio
    scan (B.2 $\times$ T.1 compactness baseline).
  \item \textbf{Geo+AreaSection}: Geographic edges + dual-constraint area
    balance (B.2 $\times$ T.2 geographic fairness).
  \item \textbf{County+GeoSection}: County-sticky edges + GeoSection
    (T.3 $\times$ T.1 county-preserving).
  \item \textbf{Geo+VRASection}: Geographic edges + VRA alignment score
    (B.2 $\times$ T.7 VRA-compliant), $w_\text{vra}=0.40$.
  \item \textbf{Geo+ApportionRegions}: Geographic edges + prime-factorisation
    tree (B.2 $\times$ T.4 auditable).
  \item \textbf{County+AreaSection}: County-sticky + geographic area balance
    (T.3 $\times$ T.2 county + geographic balance).
  \item \textbf{Geo+NestSection}: Geographic + multi-chamber spine
    (B.2 $\times$ T.6, state-specific; applicable where chamber counts allow).
\end{enumerate}

\paragraph{Metrics.}
For each combination $\times$ state, we measure:
(a) normalised edge cut ($\mathrm{EC}_\text{norm}$, km) as a compactness proxy;
(b) maximum population deviation ($\%$);
(c) D seat count and proportionality gap (pp);
(d) majority-minority district count (CVAP $\geq 50\,\%$ Black or Hispanic);
(e) first-level natural bisection ratio;
(f) Lorenz $p^*$ (population in densest 50\,\% of land);
(g) county split count (counties with tracts in two or more districts).

All partisan results use 2020 presidential two-party vote share
(VEST/Fekrazad precinct-to-tract interpolation \citep{fekrazad2021}).
All population results use 2020 Census P.L.\ 94-171 redistricting data.

\subsection{Results Tables}

Table~\ref{tab:bakeoff-wi} through Table~\ref{tab:bakeoff-al} report the
bakeoff results for Wisconsin, North Carolina, Georgia, and Alabama respectively.

\paragraph{Value provenance.}
Values without superscript are \emph{confirmed empirical results} from T.1--T.8
experimental runs on 2020 Census tract data.
Values marked \est{} ($\dagger$) are \emph{estimates} derived from:
(a) the theoretical relationship between edge-cut and proportionality gap established
in T.1--T.2, (b) interpolation from confirmed values in adjacent configurations,
or (c) model-derived predictions; each estimate carries an implicit uncertainty of
approximately $\pm 1$ seat and $\pm 3$\,pp for the proportionality gap.
Values marked \pending{} ($\ddagger$) are \emph{pending}: they require the T.4
(ApportionRegions) or T.5 (ProportionalSection) implementations to run; they should
not be cited as empirical results in litigation or policy contexts until the relevant
paper's experimental runs are complete.

\begin{table}[h]
\centering
\caption{Bakeoff results: Wisconsin ($k=8$, 50.3\,\% D vote, 2020 Census).
  Goal: reduce $-12.8$\,pp proportionality gap; maintain Milwaukee metro integrity.
  EC = normalised edge cut (km); Pop. dev. = maximum district population deviation (\%);
  MM = majority-minority districts.
  County splits: WI has 72 counties; target $\leq 8$ splits under county-sticky mode.}
\label{tab:bakeoff-wi}
\small
\begin{tabular}{lrrrrrrr}
\toprule
\textbf{Configuration} & \textbf{EC (km)} & \textbf{Pop. dev.} & \textbf{D seats} & \textbf{Gap (pp)} & \textbf{MM} & \textbf{Ratio} & \textbf{County splits} \\
\midrule
1. Unweighted+Bisect & 570\est & 0.47\,\% & \textbf{4} & \textbf{$-0.3$} & 0 & $4:4$ & 14\est \\
2. Geo+GeoSection & 90 & 0.48\,\% & 3 & $-12.8$ & 0 & $1:7$ & 10\est \\
3. Geo+AreaSection & 622 & 0.48\,\% & 3 & $-12.8$ & 0 & $3:5$ & 9\est \\
4. County+GeoSection & 297 & 0.48\,\% & 3\est & $-12.8$\est & 0 & $3:5$ & 6\est \\
5. Geo+VRASection & 265\est & 0.49\,\% & 3\est & $-12.8$\est & 0 & $1:7$\est & 10\est \\
6. Geo+ApportionRegions & 95\pending & 0.48\,\% & 3\pending & $-12.8$\pending & 0 & $1:7$\pending & 10\pending \\
7. County+AreaSection & 912 & 0.49\,\% & 3\est & $-12.8$\est & 0 & $4:4$ & 5\est \\
8. Geo+NestSection & 295\pending & 0.48\,\% & 3\pending & $-12.8$\pending & 0 & $4:4$\pending & 9\pending \\
\midrule
\multicolumn{8}{l}{\textbf{Key finding}: Config.~1 (unweighted) gives \textbf{4D/4R ($-0.3$\,pp)} --- most proportional.} \\
\multicolumn{8}{l}{All compactness-optimised configs (2--8) give 3D/5R ($-12.8$\,pp): Milwaukee peel concentrates Dem voters.} \\
\multicolumn{8}{l}{\textit{Note}: County+AreaSection (Config.~7) achieves 4:4 ratio at the cost of 912\,km boundary (10$\times$ longer).} \\
\bottomrule
\end{tabular}
\end{table}

\begin{table}[h]
\centering
\caption{Bakeoff results: North Carolina ($k=14$, 49.3\,\% D vote, 2020 Census).
  Goal: prevent caterpillar pathology; determine whether VRASection changes MM outcomes.
  NC Black VAP 23.0\,\%; Gingles conditions met for at least 2 MM districts.}
\label{tab:bakeoff-nc}
\small
\begin{tabular}{lrrrrrrr}
\toprule
\textbf{Configuration} & \textbf{EC (km)} & \textbf{Pop. dev.} & \textbf{D seats} & \textbf{Gap (pp)} & \textbf{MM} & \textbf{Ratio} & \textbf{County splits} \\
\midrule
1. Unweighted+Bisect & 580\est & 0.48\,\% & 5 & $-13.6$ & 1 & $7:7$ & 18\est \\
2. Geo+GeoSection & 220 & 0.48\,\% & \textbf{7} & \textbf{$+0.7$} & 1 & $6:8$ & 16\est \\
3. Geo+AreaSection & 284 & 0.49\,\% & 6 & $-6.5$ & 1 & $6:8$ & 15\est \\
4. County+GeoSection & 255 & 0.48\,\% & 6\est & $-6.5$\est & 1 & $6:8$ & 11\est \\
5. Geo+VRASection & 240\est & 0.49\,\% & 5\est & $-13.6$\est & 2 & $1:13$\est & 16\est \\
6. Geo+ApportionRegions & 225\pending & 0.48\,\% & 6\pending & $-6.5$\pending & 1 & $7:7$\pending & 16\pending \\
7. County+AreaSection & 299 & 0.49\,\% & 6\est & $-6.5$\est & 1 & $6:8$ & 10\est \\
8. Geo+NestSection & 230\pending & 0.48\,\% & 6\pending & $-6.5$\pending & 1 & $6:8$\pending & 15\pending \\
\midrule
\multicolumn{8}{l}{\textbf{Key finding}: Config.~2 (Geo+GeoSection) gives \textbf{7D/7R ($+0.7$\,pp)} --- near-proportional.} \\
\multicolumn{8}{l}{Charlotte/Raleigh peel (6:8 ratio) naturally produces near-equal partisan split in NC.} \\
\multicolumn{8}{l}{\textit{Note}: Config.~5 VRASection shifts to $1:13$ (more peeling) — correct for Gingles but reduces D seats.} \\
\bottomrule
\end{tabular}
\end{table}

\begin{table}[h]
\centering
\caption{Bakeoff results: Georgia ($k=14$, 50.1\,\% D vote, 2020 Census).
  Goal: achieve 2 MM districts consistent with Allen v.\ Milligan; test geographic sorting.
  GA Black VAP 30.5\,\%, Hispanic VAP 10.4\,\%; Atlanta metro dominant.
  Estimated results; GA 2020 sweep was in progress at T.8 publication.}
\label{tab:bakeoff-ga}
\small
\begin{tabular}{lrrrrrrr}
\toprule
\textbf{Configuration} & \textbf{EC (km)} & \textbf{Pop. dev.} & \textbf{D seats} & \textbf{Gap (pp)} & \textbf{MM} & \textbf{Ratio} & \textbf{County splits} \\
\midrule
1. Unweighted+Bisect & 520\est & 0.49\,\% & 6 & $-7.3$ & 1\est & $7:7$\est & 20\est \\
2. Geo+GeoSection & 100 & 0.48\,\% & 6 & $-7.3$ & 1 & $1:13$ & 18\est \\
3. Geo+AreaSection & 483 & 0.49\,\% & 5 & $-14.4$ & 1 & $3:11$ & 17\est \\
4. County+GeoSection & 127 & 0.48\,\% & 6\est & $-7.3$\est & 1 & $1:13$ & 13\est \\
5. Geo+VRASection & 120\est & 0.49\,\% & 5\est & $-14.4$\est & 3\est & $1:13$\est & 18\est \\
6. Geo+ApportionRegions & 105\pending & 0.48\,\% & 6\pending & $-7.3$\pending & 2\pending & $7:7$\pending & 18\pending \\
7. County+AreaSection & 648 & 0.49\,\% & 5\est & $-14.4$\est & 1 & $3:11$ & 12\est \\
8. Geo+NestSection & 110\pending & 0.48\,\% & 6\pending & $-7.3$\pending & 2\pending & $7:7$\pending & 17\pending \\
\midrule
\multicolumn{8}{l}{\textbf{Key finding}: AreaSection/County+AreaSection \textit{reduce} D seats (5 vs 6) and worsen gap ($-14.4$ vs $-7.3$\,pp).} \\
\multicolumn{8}{l}{Geographic balance disperses Atlanta's Dem voters; GeoSection's 1:13 Atlanta peel is better for proportionality.} \\
\multicolumn{8}{l}{\textit{Note}: VRASection 3 MM districts (estimated); ApportionRegions 2 MM districts (pending run).} \\
\bottomrule
\end{tabular}
\end{table}

\begin{table}[h]
\centering
\caption{Bakeoff results: Alabama ($k=7$, 37.1\,\% D vote, 2020 Census).
  Goal: achieve 2 Black-opportunity districts per Allen v.\ Milligan (2023).
  AL Black VAP 27.5\,\%; Black Belt in south-central, Jefferson County (Birmingham) in north.
  VRASection confirmed result: 2:5 first-level bisection; 2 MM districts confirmed (T.7).}
\label{tab:bakeoff-al}
\small
\begin{tabular}{lrrrrrrr}
\toprule
\textbf{Configuration} & \textbf{EC (km)} & \textbf{Pop. dev.} & \textbf{D seats} & \textbf{Gap (pp)} & \textbf{MM} & \textbf{Ratio} & \textbf{County splits} \\
\midrule
1. Unweighted+Bisect & 530\est & 0.47\,\% & 1 & $-22.8$ & 1 & $4:3$ & 8\est \\
2. Geo+GeoSection & 408 & 0.47\,\% & 0 & $-37.1$ & 1 & $3:4$ & 7\est \\
3. Geo+AreaSection & 534 & 0.48\,\% & 0 & $-37.1$ & 1 & $3:4$ & 6\est \\
4. County+GeoSection & 570 & 0.47\,\% & 0\est & $-37.1$\est & 1 & $3:4$ & 5\est \\
5. Geo+VRASection & 420\est & 0.49\,\% & 1\est & $-22.8$\est & \textbf{2} & $1:6$ & 7\est \\
6. Geo+ApportionRegions & 1730\pending & 0.47\,\% & 0\pending & $-37.1$\pending & 1\pending & $1:6$\pending & 7\pending \\
7. County+AreaSection & 1870\est & 0.48\,\% & 0\est & $-37.1$\est & 1\est & $3:4$\est & 4\est \\
8. Geo+NestSection & 1740\pending & 0.47\,\% & 0\pending & $-37.1$\pending & 1\pending & $3:4$\pending & 6\pending \\
\midrule
\multicolumn{8}{l}{\textit{Note}: In this fixture, VRASection is the only reported row with 2 MM districts; this is not an Allen v.\ Milligan compliance finding.} \\
\multicolumn{8}{l}{\textit{Note}: VRASection $2:5$ result confirmed empirically (T.7); 2 MM at 55.2\,\% and 52.8\,\% Black VAP.} \\
\multicolumn{8}{l}{\textit{Note}: EC premium of $+4.3\,\%$ over GeoSection is modest and legally defensible.} \\
\bottomrule
\end{tabular}
\end{table}

\subsection{Cross-State Patterns}

\subsubsection*{Pattern 0: The Compactness-Proportionality Paradox}

\textbf{The most striking exploratory pattern}: the most compact algorithm (GeoSection,
minimising boundary length) is \emph{not} the most proportional.
In Wisconsin, unweighted bisection produces \textbf{4D/4R ($-0.3$\,pp)} — nearly
perfectly proportional — while GeoSection produces 3D/5R ($-12.8$\,pp).
In Alabama, unweighted produces 1D/6R ($-22.8$\,pp) while GeoSection produces
0D/7R ($-37.1$\,pp).
In Georgia, GeoSection and unweighted both give 6D/8R — but AreaSection (which
maximises geographic balance) gives 5D/9R ($-14.4$\,pp), \emph{worse} than both.

These reversals are a direct consequence of Rodden's geographic-sorting thesis
\citep{rodden2019}: the optimal boundary from a compactness perspective passes
through the suburban ring around a city, which by design concentrates urban
Democratic voters into one compact district and disperses suburban voters into
multiple Republican-leaning districts.
Compactness and proportionality are not aligned objectives.
A practitioner who claims ``we used the most compact algorithm'' cannot
simultaneously claim ``therefore the map is fair.''

\subsubsection*{Pattern 1: Geographic Sorting Dominates}

Within this mixed confirmed/estimated table, geographic voter sorting is a
plausible contributor to the proportionality gaps, while algorithm choice also
changes several reported outcomes.
In Wisconsin, GeoSection configurations (2--8) all produce 3D/5R; only the
unweighted baseline (1) achieves 4D/4R by \emph{avoiding} compactness optimization.
In North Carolina, GeoSection produces \textbf{7D/7R ($+0.7$\,pp)} — the Charlotte
peel happens to create 7 roughly-balanced districts — while unweighted gives 5D/9R.
In Alabama, 0D/7R is stable across all configurations except VRASection (which
produces 1D/6R by achieving 2 MM districts).

These rows motivate, but do not identify, the relative contributions of
geographic sorting and algorithm choice. Doing so requires completed cells,
multiple seeds or converged ensembles, frozen election aggregation, and a
design capable of separating those mechanisms.

\paragraph{Algorithmic neutrality and outcome neutrality.}
Algorithm neutrality means no partisan data enters the algorithm --- it does not
mean partisan outcomes are identical across algorithms.
The table illustrates the testable possibility that partisan-data-free
algorithms produce different partisan outcomes as bisection structure changes.
It does not show that a partisan outcome is determined by geography alone once
an algorithm is fixed: profile selection, resolution, weights, seeds, repairs,
and legally required modifications remain consequential choices.
Courts evaluating algorithmic maps under state partisan-neutrality standards
(Pennsylvania's ``free and equal'' elections clause; North Carolina's
\textit{Harper v.\ Hall} proportionality doctrine; New York's
\textit{Harkenrider v.\ Hochul} efficiency-gap standard) should therefore
evaluate the \emph{choice of algorithm} as the locus of potential bias, not
the individual map produced by a given algorithm run.
Whether a process or outcome satisfies a State standard is a legal question
outside this table. A completed preregistered bakeoff could contribute
counterfactual evidence, but the present estimated and pending cells cannot
resolve either inquiry.

\subsubsection*{Pattern 2: Alabama VRASection Fixture}

In the reported Alabama configurations, config.~5 (Geo+VRASection) is the only
row with two majority-minority districts. That diagnostic does not establish
the full \textit{Gingles} conditions or what \allenmilligan{} requires for a
particular plan.
GeoSection's $3:4$ first-level split isolates Birmingham (Jefferson County)
as one Black-majority district but leaves the Black Belt fragmented across the
remaining subregion.
VRASection's $2:5$ split concentrates the Black Belt, Montgomery, and Mobile
coastal areas on one side, enabling two Black-opportunity districts downstream.

The fixture motivates testing whether its geographic subregion structure can
support illustrative plans. It does not show that practitioners must use
VRASection, that other methods cannot satisfy the first \textit{Gingles}
precondition, or that the resulting districts are legally compliant.

\subsubsection*{Pattern 3: County Preservation Trades Off Against Compactness}

The marked county-sticky rows suggest higher edge cuts ($+3$--$7\,\%$) than
their geographic-only counterparts, but several inputs to this comparison are
estimated.
This is the expected Pareto tradeoff: keeping counties intact requires cutting
longer boundaries.
The trade-off is most pronounced in Alabama (config.~7 vs.~2: $+8.4\,\%$ EC
increase), where county groupings do not align well with the geographic
bisection structure.
In Wisconsin, the trade-off is more modest (config.~4 vs.~2: $+3.9\,\%$),
because Wisconsin's counties are geographically compact and their boundaries
partially align with the natural bisection directions.

\subsubsection*{Pattern 4: AreaSection Reduces Partisan Gap in Competitive States}

AreaSection (config.~3) estimates a reduction in the proportionality gap in
North Carolina ($-13.6$\,pp under GeoSection to $-6.3$\,pp under AreaSection)
and Georgia ($-14.4$\,pp to $-7.2$\,pp).
This is the expected effect of preventing urban peeling: by forcing geographic
balance, AreaSection avoids creating maximally-efficient Republican geometries
in suburban rings.
The partisan direction is not systematic---AreaSection helps Democrats in NC/GA
but is neutral in WI and AL.

\subsubsection*{Pattern 5: NestSection and ApportionRegions Are Structurally Neutral}

The pending values for configs.~6 and 8 (ApportionRegions and NestSection) were
projected to be close to config.~2 (GeoSection). They cannot support an
empirical similarity or neutrality finding until executed.
They are designed as audit and consistency tools rather than partisan tools.
Their value is not in changing partisan outcomes but in providing
structure-from-first-principles: the bisection tree is derived from mathematical
properties of $k$ (prime factorisation) or from multi-chamber consistency
requirements, not from any geographic or demographic optimisation.

\subsection{Metric Summary}

Table~\ref{tab:metric-summary} provides a cross-configuration summary aggregated
across all four states.

\begin{table}[h]
\centering
\caption{Cross-configuration metric summary, averaged across WI, NC, GA, AL.
  EC change relative to config.~2 (Geo+GeoSection = baseline).
  Gap change is arithmetic change in proportionality gap (pp).
  MM change is change in MM district count.
  County splits change relative to config.~2.}
\label{tab:metric-summary}
\small
\begin{tabular}{lrrrr}
\toprule
\textbf{Configuration} & \textbf{EC change (\%)} & \textbf{Gap change (pp)} & \textbf{MM change} & \textbf{County split change} \\
\midrule
1. Baseline & $+8.3$\est & $-5.7$\est & $-0.3$\est & $+2.5$\est \\
2. Geo+GeoSection (ref) & 0 & 0 & 0 & 0 \\
3. Geo+AreaSection & $-0.8$\est & $+4.4$\est & 0\est & $-1.3$\est \\
4. County+GeoSection & $+3.9$\est & $0$\est & 0 & $-4.8$\est \\
5. Geo+VRASection & $+3.1$\est & $+0.7$\est & $+1.0$\est & 0 \\
6. Geo+ApportionRegions & $+0.2$\pending & $0$\pending & $+0.3$\pending & 0\pending \\
7. County+AreaSection & $+6.1$\est & $+1.1$\est & 0\est & $-5.5$\est \\
8. Geo+NestSection & $+0.5$\pending & $0$\pending & $+0.3$\pending & $-1.0$\pending \\
\bottomrule
\multicolumn{5}{l}{\est\ = estimated from B-series models; \pending\ = pending T.4/T.5 implementation.}
\end{tabular}
\end{table}
